\documentclass[12pt]{article}
\usepackage{hyperref}
\usepackage{amsmath, amssymb, amsfonts}
\usepackage[margin=1.2cm]{geometry}
\usepackage{xcolor}
\usepackage{graphicx}
\usepackage{enumitem, inconsolata}
\usepackage{tikz}
\usepackage{circuitikz}[scale=1.2]
\usetikzlibrary{patterns}
\parindent 0px
\newcommand{\W}{\(\Omega\)}
\newcommand{\w}{\(\omega\)}
\newcommand{\lb}{\\$\left|\rightarrow\right.$}
\newcommand{\enter}{\\\textcolor{white}{1}}
\newcommand{\sub}[1]{\textsubscript{#1}}
\newcommand{\super}[1]{\textsuperscript{#1}}
\ExplSyntaxOn
\NewDocumentCommand{\bo}{m}
 {
   \bold_commas:n { #1 }
 }

\cs_new:Npn \bold_commas:n #1
 {
   \seq_set_split:Nnn \l_tmpa_seq { , } { #1 }
   \seq_map_indexed_function:NN \l_tmpa_seq \__bold_commas_aux:nn
 }

\cs_new:Npn \__bold_commas_aux:nn #1 #2
 {
   \textbf{#2}
   \int_compare:nNnTF { #1 } < { \seq_count:N \l_tmpa_seq }
     { , }
     { }
 }

\ExplSyntaxOff

\title{Propagation and Antenna}
\author{by PUL079BEI008}
\date{\today}

\begin{document}
\maketitle
\vspace{10.5cm}
\textbf{Notes}
\begin{itemize}[noitemsep]
	\item Our subject code is EX 653.
	\item This document encorporates for both new and old question paper, both being EX 653.
	\item If a question was asked in Regular Exam, it is kept in \bo{boldened} form. If it was asked in back exam, it is in regular font.
	\item Months are marked as: 
	\begin{itemize}[noitemsep, topsep=0pt]
		\item Ba: Baisakh
		\item Jth: Jestha
		\item Asa: Ashar
		\item Shr: Shrawan
		\item Bh: Bhadra
		\item Ash: Ashwin
		\item Ka: Kartik
		\item Mng: Mangsir
		\item Po: Poush
		\item Ma: Magh
		\item Ch: Chaitra
	\end{itemize}
\end{itemize}
\pagebreak
\tableofcontents
\pagebreak

\section{Introduction}
  \begin{center} (6 Hours / 10-14 Marks) \end{center}
 
  \subsection{Retarded Potentials: EM wave generation}
    \begin{enumerate}
      \item What is retarded potential? \hfill [2] (\bo{76 Bh, 72 Ash})
        \lb and state its importance. \hfill [2] (\bo{74 Bh})

      \item What is retarded time? \hfill [2] (82 Ka, 80 Ash)

      \item Write the expression of retarded scalar potential and vector potential of infinitesimal dipole antenna in frauhofer for field region. \hfill [4] (\bo{72 Ash})
        \lb Explain retarded vector and scalar potential with mathematical expression. \hfill [5] (80 Ash)
    \end{enumerate}

    \subsubsection{Conduction current}
      \begin{enumerate}
        \item Explain how electromagnetic waves are generated by a conductor. \hfill [4] (\bo{76 Bh})
      \end{enumerate}
    \subsubsection{Short uniform current dipole}
      \begin{enumerate}
        \item Explain the operation of infinitesimal dipole with the help of mathematical relations and the field pattern. \hfill [4] (71 Ma) [6] (\bo{80 Ch, 74 Ba}) [8] (\bo{70 Bh})
      \end{enumerate}

    \subsubsection{Radiated electric and magentic fields}

  \subsection{Short uniform current dipole, short dipole and long dipole}
    \begin{enumerate}
      \item Find out the retarded vector potential of an electric current at a distance, r from a point source for the far field region. \hfill [4] (82 Ka)

      \item During the transmission mode, explain the mechanism involved by which the electric lines of force are detached from the dipole antenna to from free space waves. \hfill [4] (70 Ma)
    \end{enumerate}
    \subsection{Radiation patterns}
    \subsection{Input impedance}

  \subsection{Introduction to Antenna}
    \begin{enumerate}
      \item Define antenna. \hfill [2] (\bo{81 Ch, 80 Ch, 79 Ch, 78 Ch, 75 Bh, 70 Bh}, 75 Ba, 73 Ma)

      \item Explain the conditions that any type of conductor can radiate electromagnetic wave or not.
        \enter\hfill [4] (\bo{75 Bh})
        \lb Explain how EM waves are generated by a conductor. \hfill [4] (75 Ba)
    \end{enumerate}

  \subsection{Antenna theorems}
    \subsubsection{Reciprocity \& Superposition}
      \begin{enumerate}
        \item State the following antenna theorem:
          \lb Reciprocity. \hfill [2] (\bo{77 Ch, 75 Bh, 74 Bh}, 73 Ma, 72 Ma) [3] (\bo{79 Ch, 78 Ch})
          \lb Superposition theorem. \hfill [2] (72 Ma) [3] (\bo{76 Bh})
      \end{enumerate}

    \subsubsection{Thevenin \& Maximum power transfer Theorem}
      \begin{enumerate}
        \item Explain the following antenna theorems: Maximum Power Transfer Theorem. \hfill [2] (\bo{77 Ch}) [3] (\bo{80 Ch, 76 Bh}, 82 Ka)
        
        \item State the following theorem: Thevenin. \hfill [3] (\bo{79 Ch})

        \item Explain superposition theorem and maximum power transfer theorem with necessary diagram. \hfill [3+3] (80 Ash)
      \end{enumerate}

    \subsubsection{Compensation}
      \begin{enumerate}
        \item Write short notes on: Compensation theorem. \hfill [2] (\bo{78 Ch, 75 Bh, 74 Bh}, 73 Ma) [3] (\bo{81 Ch, 80 Ch})
      \end{enumerate}

    \subsubsection{Equality of directional patterns}
    \subsubsection{Equivalence of receiving and transmitting impedances}

\pagebreak

\section{Antenna Parameters and Arrays}
  \begin{center} (6 Hours / 4-8 Marks) \end{center}

  \subsection{Basic Antenna Parameters}
    \begin{enumerate}
      \item Explain the following antenna parameters/short notes:
        \enter a. Receiving antenna gain (G$_r$) \hfill [1.5] (76 Ba) [3] (\bo{81 Ch})
        \enter b. Free space loss (L$_{FS}$) \hfill [3] (\bo{81 Ch})
        \enter c. Directivity. \hfill [1.5] (76 Ba) [2] (70 Ma) [3] (\bo{79 Ch}) [4] (\bo{72 Ash})
        \enter d. Polarization. \hfill [2] (70 Ma) [3] (\bo{81 Ch, 79 Ch, 79 Ch, 75 Bh, 74 Bh})
        \enter e. Linear polarization. \hfill [3] (72 Ma)
        \enter f. Input impedance. \hfill [3] (\bo{79 Ch})
        \enter g. Radiation Patterns. \hfill [2] (\bo{74 Bh}) [3] (\bo{77 Ch}, 80 Ash, 76 Ba)
        \enter h. Radiation Pattern Lobes. \hfill [2] (70 Ma)
        \enter i. Beam Width. \hfill [3] (\bo{77 Ch}, 76 Ba)
        \enter j. Beam width with mathematical derivations if necessary. \hfill [2] (72 Ma)
        \enter k. Antenna efficiency. \hfill [3] (\bo{75 Bh}) [4] (\bo{72 Ash})
        \enter l. Antenna efficiency with mathematical derivations if necessary. \hfill [2] (72 Ma)
        \enter m. Radiation Efficiency. \hfill [3] (82 Ka)
        \enter n. Bandwidth. \hfill [3] (\bo{75 Bh})
        \enter o. Half power bandwidth. \hfill [2] (\bo{70 Bh})
        \enter p. Antenna gain. \hfill [2] (\bo{74 Bh})
        \enter q. Antenna gain, with mathematical derivations if necessary. \hfill [2] (72 Ma) [3] (80 Ash)
        \enter r. Power-flux density, with mathematical derivation if necessary. \hfill [3] (80 Ash)
        \enter r. Effective Isotropic Radiated Power (EIRP). \hfill [3] (82 Ka, 80 Ash, 72 Ma) [4] (\bo{77 Ch}, 73 Ma)

      \item Why antenna radiation pattern lobes are important during antenna design? \hfill [3] (71 Ma)

      \item Explain the different conditions for a simple conductor to be a good antenna. \hfill [4] (\bo{79 Ch})

      \item Sketch the current distribution and radiation pattern of center-fed vertical dipole for following lengths. \hfill [4] (76 Ba)
        \enter a. $\lambda/2$  
        \enter b. $3\lambda/2$  
        \lb for (a), (b) and (c): $2\lambda$ \hfill [6] (72 Ma)

      \item Differentiate between: $\frac{\lambda}{2}, \lambda, 3\frac{\lambda}{2}$ length dipoles in terms of their individual radiation pattern, self impedance and directivity. \hfill [6] (70 Ma)

      \item Describe the antenna gain, antenna efficiency and directivity of antenna with mathematical derivation if necessary. \hfill [6] (75 Ba)

      \item List the parameters of antenna and explain any three of them. \hfill [6] (73 Ma)
        \lb What are basic antenna parameters? Explain briefly on any 4. \hfill [4] (71 Ma)
        \lb Explain any 5 parameters of antenna. \hfill [5] (\bo{71 Bh})
    \end{enumerate}

  \subsection{Pattern multiplication}
    \subsubsection{Linear \& 2-D antenna arrays}
    \begin{enumerate}
      \item What is the importance of antenna array? \hfill [2] (76 Ba)

      \item Why we need antenna array rather than changing the length of single antenna? \hfill [2] (72 Ma)

      \item State the principle of pattern multiplication. \hfill [2] (\bo{70 Bh})

      \item Define pattern multiplication. \hfill [2] (80 Ash)

      \item What is the importance of antenna array? \hfill [2] (\bo{77 Ch})

      \item What is a linear array? \hfill [1] (\bo{73 Bh})
    \end{enumerate}

    \subsubsection{End fire}
      \begin{enumerate}
        \item Write the features of end-fire array. \hfill [4] (\bo{78 Ch})

        \item Describe the principle of End fire and broadside array. \hfill [6] (\bo{77 Ch})

        \item Draw the resultant pattern for an End Fire Array with dipoles perpendicular to the axis of array with d = $\frac{\lambda}{2}$ \& $\alpha = \pi$. \hfill [4] (70 Ma)

        \item Derive the expressions for the total field in case of two isotropic point sources with
          \enter a. an equal amplitude and opposite phase. \hfill [3+3] (72 Ma) [4+4] (\bo{81 Ch})
        \enter b. an equal amplitude and same phase.
        \lb Derive a relation for the field intensity for the array of two element isotropic radiators and current equal in magnitude but opposite in phase. Show the condition for maxima and minima direction of intensity and draw the radiation pattern. \enter\hfill [6] (\bo{75 Bh}) [6+3] (76 Ba)
        \lb Derive a relation for the field intensity for the array of two elements isotropic radiators in various conditions. \hfill [5] (\bo{71 Bh})
        
        \item Derive a mathematical relation to calculate the field intensity of an array for two element isotropic radiators. \hfill [4] (70 Ma)
        
        \item Use the principle to obtain a wave pattern for array of two short dipoles for following cases:
          \enter a. Dipoles aligned perpendicular to the array axis with d = $\lambda/2, \alpha = 0$
          \enter b. Dipoles aligned perpendicular to the array axis with d = $\lambda/2, \alpha = \pi$ \hfill [3+3] (\bo{70 Bh})

        \item Derive the expression for far field pattern of an array of two isotropic point sources having equal amplitude and quadrature phase with d = $\lambda$/4. \hfill [6] (80 Ash)

        \item Derive the expression for the total field in case of two isotropic point sources with an equal amplitude and quadrature phase and draw the radiation filed patterns when spacing between the two point sources are $\lambda$/2 and $\lambda$/4. \hfill [5] (82 Ka)

        \item Derive the relation for the field intensity of linear array of 2 isotropic radiators with equal potential and 180$^0$ out of phase. \hfill [3] (71 Ma)
      \end{enumerate}

    \subsubsection{Broadside arrays}
      \begin{enumerate}
        \item Differentiate between End fire and broadside array. \hfill [4] (\bo{73 Bh}) [6] (76 Ba)
          \lb same, but in ``write short note'' question. \hfill [4] (75 Ba)

        \item Show the condition for broad side and end fire array with necessary diagrams. \hfill [6] (\bo{71 Bh})
      \end{enumerate}

\pagebreak

\section{Antennas Classification}
  \begin{center} (10 Hours / Marks) \end{center}

  \subsection{Broad Classification}
    \subsubsection{Isotropic Antenna}
    \subsubsection{Directional Antenna}
    \subsubsection{Omnidirectional Antenna}
      \begin{enumerate}
        \item List the characteristics of onmidirectional antenna. \hfill [2] (\bo{72 Ash})

        \item Sketch the current distribution and radiation pattern of folded dipole antenna. \hfill [3] (\bo{81 Ch})

        \item Sketch the current distribution and radiation pattern of the infinitesimal dipole antenna. \hfill [3] (82 Ka)
      \end{enumerate}
    
  \subsection{Linear Dipole/Monopole}
    \subsubsection{Short Dipole}
    \subsubsection{Half-Wave Dipole}
    \subsubsection{Folded Dipole}
    \subsubsection{Monopole}
      \begin{enumerate}
        \item Write short notes on: Monopole antenna. \hfill [3] (80 Ash)
        \item Write short notes on: Marconi antenna. \hfill [4] (\bo{77 Ch, 75 Bh, 72 Ash}, 73 Ma)

        \item List the major characteristics of Marconi antenna with necessary figures. \hfill [6] (75 Ba)

        \item Explain the working principle and design of Marconni antenna. \hfill [5] (\bo{71 Bh})
      \end{enumerate}  

  \subsection{Travelling Wave Antenna}
    \begin{enumerate}
      \item What is travelling wave antenna? \hfill [1] (71 Ma) [2] (72 Ma)

      \item List the various types of wired antenna with their radiation pattern and polarization. \hfill [4] (73 Ma)

      \item What is the main difference between standing wave antenna and traveling wave antenna? \hfill [2] (71 Ma) [3] (73 Ma)
    \end{enumerate}

    \subsubsection{Single Wire}
      \begin{enumerate}
        \item Explain the radiation mechanism and conditions of a single wire antenna that can radiate electromagnetic waves or not. \hfill [4] (\bo{81 Ch})
          \lb Explain the radiation mechanism in single wire antenna. \hfill [3] (\bo{78 Ch})
          \lb How does a single wire act as an antenna? \hfill [2] (70 Ma)
          \lb Explain how single wire can radiate EM waves with necessary radiation patterns and polarization. \hfill [4] (\bo{72 Ash})
        
        \item Explain briefly radiation mechanism in single wire antenna. \hfill [5] (\bo{73 Bh})

        \item Compare resonant and non-resonant antennas. \hfill [2] (82 Ka)

        \item Explain traveling wave long wire antenna with required diagrams. \hfill [3] (\bo{79 Ch})
      \end{enumerate}

    \subsubsection{Vee Antenna}
      \begin{enumerate}
        \item Explain working principle of a V-antenna with
          \lb structure, radiation pattern. \hfill [6] (\bo{77 Ch})
          \lb construction, adv and application. \hfill [6] (\bo{76 Bh})
          \lb construction, characteristics and neat sketch. \hfill [5] (73 Ma)
          \lb construction, principle and feature. \hfill [6] (82 Ka)
          \lb construction and characteristics. \hfill [6] (80 Ash) [7] (71 Ma, 70 Ma)
          \lb characteristics and types. \hfill [5] (72 Ma)

        \item Sketch the pattern of both types of V antennas. \hfill [2] (\bo{76 Bh})

        \item Write short notes on: V antenna. \hfill [4] (76 Ba)
      \end{enumerate}

    \subsubsection{Rhombic Antenna}
      \begin{enumerate}
        \item Explain the working principle, characteristics of rhombic antenna. \hfill [4] (\bo{70 Bh})
          \lb with advantage and applications. \hfill [7] (\bo{79 Ch})
          \lb with construction. \hfill [6] (71 Ma)
          \lb with construction and necessary diagram. \hfill [8] (\bo{78 Ch, 73 Bh})
          \lb with operation. \hfill [7] (\bo{72 Ash})

        \item Write short notes on: Rhombic Antenna. \hfill [3] (\bo{76 Bh})

        \item Explain the working principle and design of Rhombic antenna. \hfill [5] (\bo{71 Bh})
      \end{enumerate}

  \subsection{Aperture Antennas}
    \begin{enumerate}
      \item Write short notes: Aperture Antenna. \hfill [3] (\bo{81 Ch})
        \lb What is aperture antenna. \hfill [1] (\bo{75 Bh})

      \item What are the advantages of aperture antenna? \hfill [2] (71 Ma)

      \item How are aperture type of antennas different from the conventional antennas? \hfill [2] (\bo{77 Ch})
    \end{enumerate}

    \subsubsection{Horn Antenna}
      \begin{enumerate}
        \item Write short notes: Horn Antenna. \hfill [3] (\bo{79 Ch}, 72 Ma)
        
        \item Write short notes: Pyramidal horn antenna. \hfill [3] (\bo{73 Bh})

        \item Explain the theory of radiation from a rectangular horn. \hfill [4] (\bo{77 Ch})

        \item Derive the relation for flare angle and length of a pyramidal horn antenna. \hfill [5] (\bo{70 Bh})

        \item List out the type of horn antenna and explain rectangular horn briefly. \hfill [6] (71 Ma)
      \end{enumerate}

    \subsubsection{Slot Antenna}

    \subsubsection{Reflector Antenna}
      \begin{enumerate}
        \item What is reflector antenna? \hfill [1] (\bo{76 Bh})

        \item Explain the special features of reflector antenna and discuss on different types of feed used with necessary diagrams. \hfill [8] (\bo{78 Ch})

        \item Write short notes on: Corner reflector antenna. \hfill [4] (76 Ba)

        \item Parabolic Reflector
        \begin{enumerate}
          \item Explain parabolic reflector:
            \lb operation. \hfill [5] (\bo{79 Ch})
            \lb with diagrams and its feeding systems. \hfill [6] (\bo{76 Bh})
            \lb construction, working principle and feeding mechanism. \hfill [6] (\bo{75 Bh}) [8] (\bo{73 Bh})

          \item Write short notes on: Parabolic reflector antenna. \hfill [4] (75 Ba)

          \item Describe the cassegrain method of feeding parabolic reflectors. \hfill [6] (73 Ma)

          \item What is Cassegrain feed system. \hfill [3] (\bo{79 Ch})
        \end{enumerate}

        \item Write short notes on: Cassegrain antenna. \hfill [3] (\bo{80 Ch}, 82 Ka)
      \end{enumerate}

  \subsection{Array Antenna}
    \subsubsection{Yagi-Uda}
      \begin{enumerate}
        \item Explain Yag-Uda antenna with construction, working principle:
          \lb with design, advantage and application. \hfill [7] (\bo{75 Bh})
          \lb with design. \hfill [5] (73 Ma) [8] (\bo{74 Bh})
        
        \item What is the main role of driving element of Yagi-Uda antenna? \hfill [3] (\bo{80 Ch})

        \item Explain significance of parasitic elements in Yagi-Uda array. \hfill [2] (80 Ash, 76 Ba, 75 Ba)

        \item Name the parasitic elements used in Yagi-Uda array. \hfill [1] (80 Ash, 76 Ba, 75 Ba)

        \item What are the advantages of Yagi antenna? \hfill [2] (72 Ma)

        \item Design a 3 element Yagi antenna for the operating frequency band of 2785 MHz to 2895 MHz and dipole as a driven element element. Take an effective length of the dipole 0.48$\lambda_c$ and consider spacing is 0.15$\lambda_c$. \hfill [6] (82 Ka)

        \item Design a 3-element Yagi antenna with the operating frequency of 3 GHz and dipole as driven element. \hfill [4] (72 Ma)

        \item Design 5 element Yagi-Uda antenna for the operating frequency band of 785 MHz to 895 MHz and dipole as a driven element. Take an effective length of the dipole 0.48 $\lambda_C$. \hfill [6] (\bo{81 Ch})

        \item Design five-element Yagi-Uda antenna with operating frequency 760 MHz and length of driver element $0.48 \lambda_C$. \hfill [6] (\bo{80 Ch})
          \lb for 100 MHz, effective length = 0.48$\lambda_c$ and spacing = 0.15$\lambda_c$. \hfill [7] (\bo{76 Bh})
          \lb for 800 MHz, effective length = 0.48$\pi$ spacing. \hfill [6] (73 Ma)
          \lb for 100 MHz, length not given. \hfill [6] (\bo{72 Ash})
      \end{enumerate}

    \subsubsection{Log-Periodic}
      \begin{enumerate}
        \item Explain the construction, working principle of Log-periodic antenna. \hfill [8] (70 Ma)
          \lb with charactersitics. \hfill [8] (70 Ma)
          \lb with characteristics and necessary diagram. \hfill [6] (\bo{81 Ch})
          \lb with design and necessary diagram. \hfill [7] (73 Ma)
          \lb with design. \hfill [8] (\bo{71 Bh}) [10] (\bo{80 Bh})

        \item Compare Yagi-Uda antenna with Log periodic dipole array. \hfill [2] (\bo{71 Bh}) [3] (80 Ash) [4] (76 Ba, 75 Ba)

        \item Write short notes on: Logarithmic antenna. \hfill [4] (\bo{74 Bh})
      \end{enumerate}

    \subsubsection{Phased Array}

    \subsection{Broadband \& Special Elements}
    \subsubsection{Helical Antenna}
      \begin{enumerate}
        \item Explain the helical antenna with respect to structural design, working principle in axial mode and normal mode, their merits, demerits and applications. \hfill [8] (\bo{77 Ch})

        \item Write short notes on: Helical Antenna. \hfill [3] (\bo{73 Bh}, 80 Ash) [4] (\bo{74 Bh})
      \end{enumerate}

    \subsubsection{Microstrip Antenna}
      \begin{enumerate}
        \item Write short notes on: Microstrip antenna. \hfill [3] (\bo{79 Ch, 78 Ch})
          \lb Write short notes on: Printed antenna. \hfill [3] (\bo{73 Bh})
      \end{enumerate}
    \subsubsection{Spiral Antenna}
    
    \subsubsection{Loop Antenna}

  \subsection{Other}
    \subsubsection{Large Plane Sheet}
    \subsubsection{Small Plane Sheet}

\pagebreak

\section{Propagation and Radio Frequency Spectrum}
  \begin{center} (6 Hours / Marks) \end{center}
  
  \begin{enumerate}
    \item Write short notes on: Radio frequency spectrum. \hfill [3] (\bo{78 Ch, 72 Ash}) [4] (\bo{75 Bh})
      \lb Define radio frequency spectrum. \hfill [2] (\bo{74 Bh})

    \item Give major propagation characteristics of VHF and UHF bands. \hfill [6] (\bo{74 Bh})
    
    \item List the main characteristics of UHF band. \hfill [4] (\bo{81 Ch})

    \item Define transmission loss. \hfill [2] (76 Ba)
  \end{enumerate}

  \subsection{Ground or Surface Wave}
    \begin{enumerate}
      \item Explain the surface wave propagation. \hfill [5] (\bo{81 Ch})
        
      \item Write major adv and disadv for surface wave propagation. \hfill [4] (\bo{81 Ch})

      \item Write the characteristics of ground wave propagation. \hfill [4] (\bo{78 Ch})

      \item Write down the factors which affect the surface wave communication. \hfill [2] (\bo{73 Bh})

      \item Write short notes on: Medium wave. \hfill [3] (80 Ash)
    \end{enumerate}
  
  \subsection{Space wave; direct and ground reflected wave, duct propagation}
    \begin{enumerate}
      \item Explain the phenomenon of Duct propagation. \hfill [4] (70 Ma) [5] (\bo{72 Ash}, 73 Ma)
        \lb in radio wave propagation. \hfill [4] (71 Ma)

      \item Write down the factors which affect the space wave communication. \hfill [5] (\bo{71 Bh})

      \item Explain the effect of space wave propagation on the ground of plane and actual earth. \hfill [5] (\bo{70 Bh}) [10] (70 Ma)
    \end{enumerate}

  \subsection{Ionospheric or sky wave; critical frequency, MUF, Skip distance}
    \begin{enumerate}
      \item What is the sky wave? \hfill [2] (80 Ash)

      \item What is MUF? \hfill [2] (\bo{75 Bh, 74 Bh})
        \lb Define MUF with necessary derivations. \hfill [2] (71 M)

      \item Describe briefly the structure of the ionosphere and its uses for radio frequency spectrum. \hfill [4] (\bo{78 Ch})
        \lb Draw the structure of ionosphere and explain the mechanism of ionosphere propagation. \hfill [6] (76 Ba)
        \lb Explain all the layers of ionosphere and their importance to radio wave communication. \hfill [6] (\bo{75 Bh})
        \lb Explai the different layers of ionosphere. \hfill [5] (71 Ma)
      
      \item Explain the operation of ionospheric communication with the help of different layers created in the atmosphere. \hfill [6] (80 Ash) [8] (\bo{77 Ch})
        \lb Explain the importance of different layers in ionospheric communication. \hfill [6] (\bo{76 Bh})

      \item Explain the ionospheric wave propagation showing its different layers. \hfill [6] (\bo{72 Ash})

      \item With a mathematical relation of refractive index of ionospheric layer derive a relation of critical frequency and MUF of radio waves with necessary explanation. Consider  the earth is not curved. \hfill [8] (\bo{71 Bh})

      \item What are the major charateristics of ionosphere? \hfill [3] (71 Ma)

      \item Define a critical frequency f$_{CRT}$. \hfill [2] (\bo{81 Ch}, 70 Ma)
        \lb Write short notes on: Critical frequency. \hfill [4] (76 Ba)
        \lb Define critical frequency with necessary derivations. \hfill [2] (71 Ma)

      \item What is skip distance? \hfill [1] (\bo{74 Bh}) [2] (72 Ma)

      \item Explain virtual height of a layer. \hfill [2] (\bo{74 Bh}, 72 Ma, 70 Ma)

      \item Define maximum usable frequency. \hfill [2] (72 Ma)

      \item Derive an expression for the relation between skip distance and maximum usable frequency (f$_{MUF}$) in SW propagation assuming that the earth is flat. \hfill [6] (\bo{81 C, 78 Ch, 74 Bh}) [7] (\bo{76 Bh}, 80 Ash)
        \lb assuming earths curvature. \hfill [7] (\bo{79 Ch}) [8] (\bo{75 Bh})
      
      \item Derive the relation between the critical frequency and maximum usable frequency. \hfill [4] (72 Ma)

      \item Briefl discuss the propagation characteristics of space wave and sky wave. \hfill [8] (\bo{73 Bh})
    \end{enumerate}

  \subsection{Troposhperic wave}
    \begin{enumerate}
      \item Find the maxium range oof tropospheric transmission for which the transmitting antenna height is 100 ft and receiving antenna is 50 ft. \hfill [5] (73 Ma)

      \item Write short notes on: Tropospheric scatter propagation. \hfill [3] (\bo{73 Bh})
    \end{enumerate}

  \subsection{RF spectrum and its propagation characteristics}
    \begin{enumerate}
      \item What is radio frequency spectrum? \hfill [2] (\bo{80 Ch})

      \item Explain characteristics of MW radio propagation. \hfill [3] (\bo{73 Bh, 71 Bh}) [4] (\bo{80 Ch})
      
      \item Write short notes on: MW propagation. \hfill [4] (73 Ma)

      \item Explain characteristics of SW radio propagation. \hfill [3] (\bo{73 Bh, 71 Bh}) [4] (\bo{80 Ch})

      \item With a neat diagram, explain the designation of radio waves according to the path they follow during propagation. \hfill [4] (\bo{70 Bh})

      \item Compare the propagation characteristics for different radio bands. \hfill [6] (\bo{70 Bh})
    \end{enumerate}

  \subsection{Numerical}
    \begin{enumerate}
      \item Find the received power in dB at a distance of 50 km over a free space with 10 GHz frequency consisting of transmitting antenna with 25 dB gain and a receiving gain of 20 dB. The power radiated by transmitted antenna is 150 dB. \hfill [6] (\bo{79 Ch})

      \item A free space LoS microwave link operating at 10 GHz consists of a transmit and receive antenna each having gain of 25 dB. Calculate the path loss of the link nd received power of 100 km distance. Derive Friis transmission equation. \hfill [3] (\bo{78 Ch})

      \item A parabolic reflector antenna having antenna efficiency 75\% is designed for 3 GHz resonant frequency with 2.5 dB waveguide loss. Find out the antenna diameter if effective isotropic radiated power (EIRP) is calculated 46 dBW and transmitting power is 500 W. \hfill [7] (\bo{77 Ch})
        \lb resonant frequency = 5 GHz, 3 dB loss, EIRP = 50 dBW, power = 600W. \hfill [8] (\bo{74 Bh})

      \item Microwave link is assumed to be free space condition. The antenna gains are each 40 dB, the frequency is 10 GHz and the path length is 90 km. Calculate the transmission path loss and received power of 10 kW. \hfill [6] (\bo{75 Bh})

      \item Find the received power (in dBm) at a distance of 0.5 km over a free space 1 GHz circuit consisting of a transmitting antenna with 25 dB gain and a receiving antenna gain of 20 dB. The power radiated by the transmitting antenna is 150 W. \hfill [5] (75 Ba)

      \item A parabolic reflector antenna having the antenna efficiency 85\% is designed for 3 GHz resonant frequency with 2.5 dB wavelength loss. Find out the antenna diameter if EIRP is calculated 46 dBW and trasmitting power is 500 W. \hfill [8] (\bo{73 Bh, 71 Bh})

      \item Assume the reflection take place at a height of 400 km and that the maximum density in the ionosphere corresponds to a 0.8 refractive index at 15 MHz. What will be the range (assume flat earth) for which the MUF is 20 MHz. \hfill [5] (70 Ma)
    \end{enumerate}

\pagebreak

\section{Propagation between Antennas}
  \begin{center} (6 Hours / Marks) \end{center}

  \subsection{Free Space propagation}
    \subsubsection{Power density of receiving antenna}
      \begin{enumerate}
        \item What is free space propagation? \hfill [2] (80 Ash)

        \item Derive Friis transmission equation for free space propagation. \hfill [7] (\bo{80 Ch}, 76 Ba)
          \lb How do you get Friis transmission equation and path loss in case of free space wave propagation? \hfill [4+2] (80 Ash) [4+3] (\bo{71 Bh})

        \item Derive Friis transmission equation. \hfill [2] (\bo{76 Bh})

        \item Write short notes on: Friis transmission equation \hfill [3] (\bo{73 Bh})

        \item Define free space communication. \hfill [2] (75 Ba)
      \end{enumerate}

    \subsubsection{Path loss}
      \begin{enumerate}
        \item Derive path loss in case of free space propagation. \hfill [2] (\bo{76 Bh}) [5] (71 Ma)
          \lb Derive complete equation including path loss using friis space communication. \hfill [8] (75 Ba)
          \lb Derive the expression for the path loss in case of radio wave propagation. \hfill [6] (\bo{73 Bh})

        \item Write short notes on: Friis equation and path loss in case of free space propagation. \hfill [3] (\bo{72 Ash})

        \item Find out distance of line of sight as: d = 3.57$\left( \sqrt{h_t} + \sqrt{h_r} \right)$ km. \hfill [3] (\bo{76 Bh})

        \item Calculate free space path loss for 6 GHz frequency, if transmitter and receiver are 10,500 km line of sight apart. \hfill [3] (71 Ma)
      \end{enumerate}

  \subsection{Plane earth propagation}
    \subsubsection{Ground Reflection}
    \subsubsection{Effective antenna heights}
    \begin{enumerate}
      \item Write short notes on: Effective antenna height. \hfill [3] (\bo{78 Ch})
    \end{enumerate}
    
    \subsubsection{The Two Ray}
  \subsection{Propagation Model, Path Loss}
  \subsection{Fresnel Zones and Knife edge diffraction}
    \begin{enumerate}
      \item Write short notes on Fresnel knife edge diffraction. \hfill [2] (71 Ma, 70 Ma) [3] (\bo{81 Ch, 80 Ch, 79 Ch, 77 Ch, 76 Bh, 72 Ash}, 72 Ma) [4] (\bo{75 Bh, 74 Bh}, 75 Ba, 73 Ma) [5] (\bo{70 Bh})

      \item Write short notes on Fresnel Zones. \hfill [2] (71 Ma) [3] (\bo{78 Ch})\\
        (\textit{Original ``Firesnel" was assumed to be Fresnel})
    \end{enumerate}

  \subsection{Numerical}
    \begin{enumerate}
      \item Find out the  free space loss (L$_{FS}$) equation that equals to 22 + 20 $\log(\frac{d}{\lambda})$, where d is the distance between T$_X$ and R$_X$ antennas and $\lambda$ is the wavelength of the frequency being used. \hfill [3] (\bo{80 Ch})

      \item What is the maximum power that can be received over a distance of 15.5 km line-of-sight free space with a 3.4 GHz frequency consisting of the transmitting antenna gain of 25 dB and receiving antenna diameter 6.4m with 75\% antenna efficiency. Where transmitted power is 250W. \hfill [7] (76 Ba)

      \item The antenna of a TV transmitter is located at a height of 125m above ground level. Calculate the distance upto which the LOS communication is possible if the height of receiving antenna is to be 9m. \hfill [6] (\bo{74 Bh})

      \item Find out the line of sight of distance between the transmitting antenna and receiving antenna. If the transmitting antenna height is 45 m, and the receiving antenna height is 25 m. [Given the radius of Earth is 6,378 km]. \hfill [6] (\bo{72 Ash})

      \item A 100 MHz circuit consists of a transmitting and receiving antenna of 30 dB and 25 dB gains respectively. The power radiated by the transmitting antenna is 120 W. Using Friis transmission equation, find the received power at a distance of 0.75 km over a free space. \hfill [6] (\bo{72 Ash})
    \end{enumerate}

\pagebreak

\section{Optical fibres}
  \begin{center} (11 Hours / Marks) \end{center}
  
  \subsection{Optical Fibre Communication System}
    \begin{enumerate}
      \item What is optical fiber? \hfill [1] (72 Ma)

      \item Explain the block diagram of optical communication system. \hfill [3] (\bo{78 Ch}, 71 Ma)

      \item Draw the optical fiber communication system. \hfill [3] (\bo{74 Bh})

      \item What are the various elements of an optical communication system? Explain each element in brief. \hfill [6] (76 Ash) [8] (\bo{73 Bh, 72 Ash})

      \item What properties do you think to be considered when an optical fiber is to be selected for any communication system? \hfill [6] (72 Ma)
    \end{enumerate}

    \subsubsection{Merits/Demerits}
      \begin{enumerate}
        \item What are the advantages and disadvantages of optical fiber over wireless communication system? \hfill [5] (\bo{81 Ch})  
          \lb State adv of optical fiber compared to microwave communication link. \hfill [4] (\bo{78 Ch})
          \lb Explain adv and disadv of optical fibers over metal wire communication. \hfill [5] (\bo{76 Bh}) [6] (\bo{74 Bh})
          \lb Briefly explain the adv of optical fiber over metalled wire. \hfill [6] (70 Ma/de)
          \lb List out the adv and disadv of optical communication over cable communication. \hfill [5] (71 Ma) 

        \item Explain the block diagram of optical communication system. \hfill [3] (\bo{78 Ch})

        \item Where are optical fibers most widely used? \hfill [2] (\bo{76 Bh, 70 Bh})

        \item Compare optical fiber communication with cable and radio communication system? Explain each element in brief. \hfill [8] (\bo{73 Bh})
      \end{enumerate}

  \subsection{Types of optical fibre and its Structural differences}
    \begin{enumerate}
      \item Explain the contribution of microscope and macroscopic fiber bends towards the bending losses in optical fiber. \hfill [7] (\bo{77 Ch})

      \item Explain the construction and propagation mechanism and application of different types of optical fibers. \hfill [5] (\bo{76 Bh}) [8] (73 Ma) [10] (\bo{71 Bh})

      \item Explain the construction and types of optical fiber with necessary diagrams. \hfill [5] (76 Ba)
        \lb Compare the configuration of different types of fiber. \hfill [4] (\bo{72 Ash})
      
      \item What are the various feature of graded index fiber? \hfill [3] (76 Ba)

      \item Write short notes on: Multimode graded index fiber. \hfill [4] (73 Ma)

      \item Explain the refractive index profile and ray transmission in a multimode graded index fiber. \hfill [5] (76 Ba)

      \item What are the advantages and disadvantages of multimode fiber over the single mode fiber? \hfill [3] (70 Ma) [5] (72 Ma)
        \lb Mention the advantages of single mode fiber. \hfill [2] (\bo{72 Ash})
    \end{enumerate}

  \subsection{Light propagation characteristics}
    \begin{enumerate}
      \item Explain the working principle of optical fiber communication system. \hfill [3] (\bo{81 Ch})

      \item Write short notes on: Super refraction. \hfill [4] (\bo{74 Bh})

      \item Explain the principle of light propagation through a fiber. \hfill [2] (\bo{79 Ch})
        \lb with necessary diagram. \hfill [5] (\bo{80 Ch})
        \lb Describe with neat diagram the basic principle of total internal reflection that enables the fiber to work as a ``light conduit". \hfill [6] (\bo{75 Bh})

      \item What is an acceptance angle? \hfill [2] (\bo{81 Ch, 78 Ch, 70 Bh})

      \item Derive an expression to calculate the acceptance angle. \hfill [7] (\bo{70 Bh})
    \end{enumerate}

  \subsection{NA in optical fibre}
    \begin{enumerate}
      \item Derive the expression for numerical aperture of a stepped index optical fiber. \hfill [5] (\bo{81 Ch})
        \lb Derive the relationship between numerical aperture and acceptance angle. \hfill [6] (\bo{80 Ch})
        \lb Derive the expression for NA with neat diagram. \hfill [5] (\bo{78 Ch})
        \lb Derive expression of NA of a stepped index optical fiber. \hfill [7] (70 Ma)

      \item Define numerical aperture. \hfill [2] (\bo{77 Ch})
        \lb with necessary figures and mathematical expression. \hfill [2] (\bo{79 Ch})
      
      \item Define acceptance angle with necessary figures and mathematical expression. \hfill [2] (\bo{79 Ch})

      \item Write short notes on: NA. \hfill [3] (71 Ma) [4] (\bo{74 Bh})
        \lb Write short notes on: NA for meridional rays in optical fiber. \hfill [3] (\bo{72 Ash})

      \item Describe the following term: Dispersion. \hfill [3] (71 Ma)
    \end{enumerate}

  \subsection{Losses}
  \begin{enumerate}
    \item Write short notes on: Optical fiber loss. \hfill [3] (\bo{79 Ch})
      
    \item Discuss loss or signal attenuation in an optical fiber with respect to absorption, scattering and bending losses. \hfill [8] (75 Ba)

    \item What are the different types of losses you can visualize in optical fiber communication? \hfill [4] (72 Ma)
  \end{enumerate}

  \subsection{Light Source and Photo Detector}
    \begin{enumerate}
      \item Write short notes on: Light source and photo detector. \hfill [3] (\bo{80 Ch})
        \lb on: Optical source and Optical detector. \hfill [3] (80 Ash) [4] (\bo{77 Ch, 75 Bh}, 75 Ba)

      \item Write short notes on: Critical frequency. \hfill [4] (76 Ba)

      \item List out different optical sources and explain the losses in optical fiber briefly. \hfill [6] (\bo{77 Ch}) [2+5] (\bo{75 Bh})
    \end{enumerate}

\section{Don't know where these are}
  \begin{enumerate}
    \item Write short notes on: Power flux density (PFD). \hfill [3] (\bo{80 Ch})
    
    \item Describe multiple HOP and explain how atmospheric ducts can be used for microwave propagation. \hfill [2+6] (\bo{79 Ch})

    \item Write short notes on: MW Propagation. \hfill [3] (\bo{76 Bh})

    \item Write short notes on: Splicing. \hfill [3] (\bo{79 Ch})
  \end{enumerate}
\end{document}
